%------------------------------------------------------------------------------
% GENERATED FILE -- DO NOT EDIT DIRECTLY.
% Source of truth: data/profile.yaml
% Regenerate with: python3 scripts/cv.py build
% Template derived from Jake Gutierrez / sb2nov (MIT License).
%------------------------------------------------------------------------------

\documentclass[letterpaper,10pt]{article}

\usepackage[T1]{fontenc}
\usepackage[utf8]{inputenc}
\usepackage{lmodern}
\usepackage{latexsym}
\usepackage[empty]{fullpage}
\usepackage{titlesec}
\usepackage{marvosym}
\usepackage[usenames,dvipsnames]{color}
\usepackage{verbatim}
\usepackage{enumitem}
\usepackage[hidelinks]{hyperref}
\usepackage{fancyhdr}
\usepackage[french]{babel}
\frenchsetup{StandardLists=true}
\usepackage{tabularx}
\usepackage{fontawesome5}
\input{glyphtounicode}

\pagestyle{fancy}
\fancyhf{}
\fancyfoot{}
\renewcommand{\headrulewidth}{0pt}
\renewcommand{\footrulewidth}{0pt}

% Margins
\addtolength{\oddsidemargin}{-0.5in}
\addtolength{\evensidemargin}{-0.5in}
\addtolength{\textwidth}{1.0in}
\addtolength{\topmargin}{-0.55in}
\addtolength{\textheight}{1.1in}

\urlstyle{same}

\raggedbottom
\raggedright
\setlength{\tabcolsep}{0in}

% Section formatting
\titleformat{\section}{
  \vspace{-6pt}\scshape\raggedright\large
}{}{0em}{}[\color{black}\titlerule \vspace{-5pt}]

% Machine-readable / ATS-parsable PDF
\pdfgentounicode=1

% Custom commands
\newcommand{\resumeItem}[1]{
  \item\small{
    {#1\vspace{-2pt}}
  }
}

\newcommand{\resumeSubheading}[4]{
  \vspace{-3pt}\item
    \begin{tabular*}{0.97\textwidth}[t]{l@{\extracolsep{\fill}}r}
      \textbf{#1} & #2 \\
      \textit{\small#3} & \textit{\small #4} \\
    \end{tabular*}\vspace{-7pt}
}

\newcommand{\resumeProjectHeading}[2]{
    \item
    \begin{tabular*}{0.97\textwidth}{l@{\extracolsep{\fill}}r}
      \small#1 & #2 \\
    \end{tabular*}\vspace{-7pt}
}

\renewcommand\labelitemii{$\vcenter{\hbox{\tiny$\bullet$}}$}

\newcommand{\resumeSubHeadingListStart}{\begin{itemize}[leftmargin=0.15in, label={}]}
\newcommand{\resumeSubHeadingListEnd}{\end{itemize}}
\newcommand{\resumeItemListStart}{\begin{itemize}[itemsep=-1pt]}
\newcommand{\resumeItemListEnd}{\end{itemize}\vspace{-5pt}}

\begin{document}

%----------HEADING----------
\begin{center}
    \textbf{\Huge \scshape Ilias Laoukili} \\ \vspace{1pt}
    \small Étudiant à la maîtrise recherche (M.Sc.) à l'Université Laval \\ \vspace{2pt}
    \href{mailto:ilias.laoukili@ulaval.ca}{\color{MidnightBlue}\faEnvelope\ ilias.laoukili@ulaval.ca} $|$
    \href{https://github.com/ilias-laoukili}{\color{MidnightBlue}\faGithub\ github.com/ilias-laoukili} $|$
    \href{https://www.linkedin.com/in/ilias-laoukili}{\color{MidnightBlue}\faLinkedin\ in/ilias-laoukili} $|$
    \href{https://ilias.laoukili.com}{\color{MidnightBlue}\faGlobe\ ilias.laoukili.com}
\end{center}

\section{Résumé de profil}
\small{
    Élève-ingénieur spécialisé en apprentissage automatique et cybersécurité, avec une expérience de recherche en réseaux de neurones sur graphes (GNN) et systèmes ML. Actuellement inscrit à la maîtrise recherche (M.Sc.) à l'Université Laval sous la direction du Prof. Talal Halabi, et conseiller en IA (emploi étudiant) au ministère de la Cybersécurité et du Numérique (MCN).
}
\vspace{-4pt}

\section{Formation}
\resumeSubHeadingListStart
  \resumeSubheading
    {Université Laval}{Québec, QC, Canada}
    {Maîtrise en informatique et IA avec mémoire (Double diplôme)}{Automne 2026 -- Présent}
    \resumeItemListStart
      \resumeItem{Maîtrise recherche avec mémoire, sous la direction du Prof. Talal Halabi.}
      \resumeItem{Axes : apprentissage automatique appliqué, cybersécurité et systèmes ML.}
    \resumeItemListEnd

  \resumeSubheading
    {ESIEE Paris}{Noisy-le-Grand, France}
    {Diplôme d'ingénieur en Intelligence Artificielle et Cybersécurité (Grande École)}{2024 -- 2027}
    \resumeItemListStart
      \resumeItem{\textbf{Moyenne actuelle : 16,08/20} (tête de promotion).}
      \resumeItem{\textbf{Notes clés :} Introduction à la cybersécurité (19,6/20), Sécurité logicielle (19,5/20), Deep Learning (19,0/20), Algorithmique des graphes (18,0/20).}
    \resumeItemListEnd

  \resumeSubheading
    {Lycée Gustave Monod}{Enghien-les-Bains, France}
    {Classes Préparatoires aux Grandes Écoles (CPGE)}{2021 -- 2024}
    \resumeItemListStart
      \resumeItem{Formation intensive de premier cycle en mathématiques avancées et physique.}
    \resumeItemListEnd

\resumeSubHeadingListEnd

\section{Expérience en recherche et ingénierie appliquée}
\resumeSubHeadingListStart
  \resumeSubheading
    {Conseiller en intelligence artificielle (Emploi étudiant)}{Nov. 2026 -- Présent}
    {Ministère de la Cybersécurité et du Numérique (MCN) -- Bureau de la transformation de l'IA}{Québec, QC, Canada}
    \resumeItemListStart
      \resumeItem{Contribution à la capacité ministérielle de réalisation et d'industrialisation de services numériques IA sécurisés et réutilisables pour le secteur public.}
      \resumeItem{Soutien à la documentation d'architecture, la préparation des données, le prototypage, les essais et l'intégration de solutions IA avec des spécialistes en sécurité et MLOps.}
    \resumeItemListEnd

  \resumeSubheading
    {Stagiaire ingénieur R\&D (IA appliquée et systèmes)}{Mai 2026 -- Août 2026}
    {Nanyang Polytechnic (NYP)}{Singapour}
    \resumeItemListStart
      \resumeItem{Développement d'un framework RAG pour les simulations énergétiques de bâtiments durables.}
      \resumeItem{Passerelle entre traitement du langage naturel et systèmes d'ingénierie en rendant les architectures \textbf{EnergyPlus} interrogeables en langage naturel.}
    \resumeItemListEnd

  \resumeSubheading
    {Tremplin Recherche -- Initiation à la recherche (Systèmes ML et GNN)}{Nov. 2025 -- Mai 2026}
    {Université Gustave Eiffel / LAMA}{Champs-sur-Marne, France}
    \resumeItemListStart
      \resumeItem{Étude de la sparsification déterministe de graphes pour optimiser l'entraînement et l'efficacité de calcul des GNN.}
      \resumeItem{Développement d'une bibliothèque en \textbf{Rust} (\href{https://github.com/ilias-laoukili/plekton}{\color{MidnightBlue}plekton}) avec bindings Python zero-copy pour générer des graphes contrôlés et évaluer la sparsification spectrale.}
      \resumeItem{Écriture de scripts \textbf{Bash} orchestrant les balayages d'hyperparamètres (parallélisation des essais légers, isolation des expériences gourmandes en mémoire pour éviter les erreurs OOM).}
    \resumeItemListEnd

\resumeSubHeadingListEnd

\section{Projets techniques}
\resumeSubHeadingListStart
  \resumeProjectHeading
    {\textbf{\href{https://github.com/ilias-laoukili/manis}{\color{MidnightBlue}manis -- Audit de durcissement Linux hors ligne}} $|$ \emph{Rust, Linux, CIS/ANSSI, SARIF}}{2026}
    \resumeItemListStart
      \resumeItem{Binaire Rust statique sûr (\textasciitilde{}1,9 Mo, zéro sous-processus) auditant images et rootfs Linux hors ligne sans droits root.}
      \resumeItem{103 règles (CIS Ubuntu 24.04, ANSSI-BP-028, NIST ; famille Debian et \texttt{linux-generic}), SARIF 2.1.0 et \texttt{diff} testés en CI sur 5 distros.}
    \resumeItemListEnd

  \resumeProjectHeading
    {\textbf{Sécurité logicielle appliquée et exploitation de vulnérabilités} $|$ \emph{ESIEE Paris}}{}
    \resumeItemListStart
      \resumeItem{\textbf{Rétro-ingénierie :} Analyse bas niveau de binaires Linux x86-64 strippés avec GDB/objdump ; contournement de protections anti-débogage (\texttt{ptrace}) et d'intégrité par patch binaire, et écriture de générateurs de clés automatisés en C.}
      \resumeItem{\textbf{Sécurité Web :} Automatisation de l'extraction de données par injection SQL aveugle (Blind SQLi) via un script d'exploitation Python dédié.}
    \resumeItemListEnd

  \resumeProjectHeading
    {\textbf{\href{https://github.com/ilias-laoukili/sentim-app}{\color{MidnightBlue}Sentim -- Deep Learning pour la reconnaissance d'émotions audio}} $|$ \emph{Python, Librosa, PyTorch, Streamlit}}{Nov. 2025}
    \resumeItemListStart
      \resumeItem{Pipeline complet de reconnaissance d'émotions audio, du prétraitement du signal brut jusqu'à l'entraînement et au déploiement.}
      \resumeItem{Conversion en spectrogrammes 2D et évaluation comparative (CNN 2 couches sur mesure et ResNet affiné surpassant la baseline Random Forest), avec tableau de bord Streamlit temps réel.}
    \resumeItemListEnd

\resumeSubHeadingListEnd

\section{Compétences techniques}
\begin{itemize}[leftmargin=0.15in, label={}]
  \small{\item{
    \textbf{Systèmes \& Bas niveau}{: Rust, PyO3, C, Bash (exécution maîtrisée en mémoire), Linux, calcul sous contrainte mémoire} \\
    \textbf{IA appliquée \& ML}{: Python, PyTorch, PyTorch Geometric, CNN/ResNet, GNN, architectures RAG/LLM, Scikit-learn} \\
    \textbf{Cybersécurité}{: Rétro-ingénierie (GDB, objdump), exploitation de vulnérabilités, Wireshark, audit de sécurité logicielle/système} \\
    \textbf{Outils \& DevOps}{: Git, Docker, CI/CD, MongoDB}
  }}
\end{itemize}

\end{document}
